% !TEX root = ../notes.tex

\documentclass[../notes.tex]{subfiles}

\begin{document}

\section{April 13}
Today we classify connected compact Lie groups.

\subsection{Classification of Compact Connected Lie Groups}
We will show the following on the homework.
\begin{proposition} \label{prop:compact-to-reductive}
	Let $K$ be a compact Lie group. Then $\op{Lie}K_\CC$ is reductive.
\end{proposition}
\begin{proof}[Sketch]
	Using integration, one can find an invariant non-degenerate bilinear form on $\op{Lie}K$, which one can use to show that the adjoint representation is semisimple. It follows that $\op{Lie}K_\CC$ can be decomposed into a sum of simple Lie algebras.
\end{proof}
For the rest of the subsection, $K^c$ will be a connected compact Lie group, and its Lie algebra will be denoted $\mf k^c$, where $\mf k\coloneqq\mf k^c_\CC$. We recall that the short exact sequence
\[0\to\mf z(\mf k)\to\mf k\to\mf k_{\mathrm{ss}}\to0\]
splits. Thus, $K^c$ will basically be a product of $G_{\mathrm{ad}}^c$ for the semisimple Lie group corresponding to $\mf k_{\mathrm{ss}}$ and some $S^1$s coming from $\mf z(\mf k)$. With this in mind, our classification will be done in two steps.
\begin{itemize}
	\item We will show that the universal cover $G^c$ of $G_{\mathrm{ad}}^c$ is compact.
	\item Then we will show that $K$ can be realized as a quotient of $G^c\times Z$ by a finite central subgroup, where $Z$ is some finite product of $S^1$s.
\end{itemize}
Let's begin with the first point. This will require a finer study of the fundamental group, which comes from the following construction.
\begin{definition}
	Let $\mf g$ be the complex semisimple Lie algebra, and let $\pi\colon G\to G_{\mathrm{ad}}$ be the universal cover of $G_{\mathrm{ad}}$. For each dominant integral weight $\lambda$, we define $\chi_\lambda$ to be the character $\ker\pi\to\CC^\times$ giving the action on $L_\lambda$. Note that $L_\lambda$ lifts to a representation of $G$ because $G$ is simply connected.
\end{definition}
\begin{remark} \label{rem:pi1-fg}
	We quickly recall why $G$ is actually a Lie group: the point is that $\pi_1(G)$ is a finitely generated group and hence countable. This holds by finding a countable locally finite cover $G$ with small balls, which shows that $G$ has the same fundamental group as a countable locally finite graph. The fundamental group of such a graph is then finitely generated.
\end{remark}
\begin{remark}
	Note that $\ker\pi$ is a discrete central subgroup: it is certainly discrete and normal (because it is a kernel). Then it is central because the adjoint action of the total group must act smoothly, so it cannot move around the discrete points.
\end{remark}
\begin{remark}
	Note that $\chi_\lambda$ is well-defined by Schur's lemma.
\end{remark}
\begin{remark}
	Because $L_{\lambda+\mu}\subseteq L_\lambda\otimes L_\mu$ (indeed, the latter has a vector of highest weight $\lambda+\mu$), we find that $\chi_{\lambda+\mu}=\chi_\lambda\chi_\mu$. Thus, $\chi_\bullet$ is additive on the dominant integral weights $P_+$, so we may extend $\chi_\bullet$ to a character
	\[\chi_\bullet\colon P\to\op{Hom}(\ker\pi,\CC^\times).\]
\end{remark}
\begin{remark}
	Because the adjoint representation $\mf g$ descends to $G_{\mathrm{ad}}\subseteq\op{Aut}(\mf g)$ for free, we conclude that $\chi_\theta=1$, where $\theta$ is the maximal root.
\end{remark}
\begin{proposition}
	Let $\mf g$ be the complex semisimple Lie algebra, and let $\pi\colon G\to G_{\mathrm{ad}}$ be the universal cover of $G_{\mathrm{ad}}$. Then the highest weight representation $L_\lambda$ descends to $G_{\mathrm{ad}}$ if and only if $\lambda\in Q$.
\end{proposition}
\begin{proof}
	% Because $G$ is semisimple, the differential provides an equivalence
	% \[\op{Rep}_\CC G\to\op{Rep}_\CC\mf g,\]
	% where these representation categories only see finite-dimensional representations. For example, this immediately implies that the former category is semisimple with irreducibles given by the highest weight representations $L_\lambda$.
	% Now, $L_\lambda$ descends to $G_{\mathrm{ad}}$ if and only if $Z$ acts trivially on $L_\lambda$, so let's measure this: because $L_\lambda$ is some irreducible representation of $G$, Schur's lemma shows that $Z$ acts by some scalars on each $L_\lambda$, so we let $\chi_\lambda\colon Z\to\CC^\times$ be the induced character. Here are some checks on this character.
	% \begin{itemize}
	% 	\item 
	% 	\item 
	% 	\item 
	% \end{itemize}
	We show the two implications separately.
	\begin{itemize}
		\item Suppose $\lambda\in Q$. On the homework, we showed that $L_{\lambda+\alpha}\subseteq\mf g\otimes L_\lambda$ for $\lambda$ large enough (i.e., provided that $(\lambda,\alpha_i^\lor)$ is large enough for $i$ large enough). Thus, $\chi_{\lambda+\alpha}=\chi_\lambda$ for $\lambda$ large enough, so $\chi_\alpha=1$ follows. (For example, this recovers the fact that $\chi_\theta=1$.) Thus, $\chi_\bullet$ descends to a character on $P/Q$. Now, $\chi\colon P/Q\to\op{Hom}(Z,\CC^\times)$ may alternatively be viewed as a character $Z\to(P/Q)^\lor$ by duality, which is also $P^\lor/Q^\lor$.
		\item For the converse, we see that $G_{\mathrm{ad}}$ admits the faithful self-dual representation $\mf g$, so $L_\lambda$ descending to $G_{\mathrm{ad}}$ implies that it descends to $G_{\mathrm{ad}}^c$, which in turn is $\otimes$-generated by $\mf g$. However, the weights of $\mf g$ are in $Q$, so it follows by taking the character that the weights of $L_\lambda$ are in $Q$, so $\lambda\in Q$.
		\qedhere
	\end{itemize}
\end{proof}
\begin{remark}
	It follows from the proof that the character $\chi\colon P/Q\to\op{Hom}(Z,\CC^\times)$ is injective.
\end{remark}
All the same arguments hold for the universal cover $G^c\to G_{\mathrm{ad}}^c$. We are now ready to calculate $\pi_1$.
\begin{theorem} \label{thm:compute-pi1-compact-lie-group}
	Let $\mf g$ be the complex semisimple Lie algebra, and let $\pi\colon G^c\to G_{\mathrm{ad}}^c$ be the universal cover of $G_{\mathrm{ad}}$, and let its kernel be $Z^c$. Then
	\[\chi\colon P/Q\to\op{Hom}(Z^c,\CC^\times)\]
	is an isomorphism.
\end{theorem}
\begin{proof}
	By \Cref{rem:pi1-fg}, we find that $Z^c$ is a finitely generated abelian group. For a choice of finite-index subgroup $Z'$, we note that $G/Z'$ is a finite cover of $G_{\mathrm{ad}}^c$ and therefore a compact Lie group which is a central extension of $G_{\mathrm{ad}}^c$ by the center $Z^c/Z'$.
	
	Notably, the finite-dimensional representations of $G^c/Z'$ are given by the sum $L_\lambda$s where $\chi_\lambda|_{Z'}$ vanishes; notably, these $\lambda$s include $Q$. We then let $P'_+\subseteq P_+$ be the subset of those dominant integral weights, and we let $P'$ be the induced subgroup. Now, $\chi_\bullet$ descends to an injective homomorphism
	\[P'/Q\to\to\op{Hom}(Z^c/Z',\CC^\times).\]
	Now, $G'$ is compact, so the Peter--Weyl theorem tells us that all characters are in fact achieved, so the $\chi_\lambda$ span a dense subset in $C(G)$, which gives the surjectivity of this map.
\end{proof}
We can collect all of our work into the following two results.
\begin{corollary} \label{cor:classify-semisimple-compact}
	Let $\mf g$ be the complex semisimple Lie algebra.
	\begin{listalph}
		\item Then the simply connected Lie group $G^c$ with Lie algebra $\mf g^c$ is compact with center $P^\lor/Q^\lor$. Furthermore, $\pi_1(G_{\mathrm{ad}}^c)\cong P^\lor/Q^\lor$.
		\item If $\Gamma\subseteq P^\lor/Q^\lor$ is a subgroup, then the category $\op{Rep}G^c/\Gamma$ is semisimple and generated by the highest weight representations $L_\lambda$ such that $\chi_\lambda|_\Gamma=1$.
		\item Decompose $\mf g$ into a sum $\bigoplus_{i=1}^n\mf g_i$ of simple Lie algebras. Then any connected Lie group $G$ with Lie algebra $\mf g^c$ is compact and takes the form
		\[\Bigg(\prod_{i=1}^nG_{i,\mathrm{ad}}^c\Bigg)/Z^c,\]
		where $Z^c=\pi_1(G^c)$. Every semisimple compact connected Lie group is of this form.
	\end{listalph}
\end{corollary}
\begin{proof}
	Here, (a) and (b) from \Cref{thm:compute-pi1-compact-lie-group}. It only remains to show the classification parts of (c), which we do in two steps.
	\begin{itemize}
		\item Note that any semisimple compact connected Lie group $G$ will have complex semisimple Lie algebra $\mf g_\CC$. Then $\mf g$ must have negative-definite Killing form in order for $G$ to be compact (otherwise, the adjoint group needs to act transitively on a form which is neither positive-definite nor negative-definite!), so $\mf g=\mf g_\CC^c$ by our classification of real forms!
		\item Once $G$ has Lie algebra $\mf g^c$, it surjects onto $G_{\mathrm{ad}}$, and this must be a covering map, so $G$ is a quotient of the universal cover $G^c$. This $G^c$ can be computed on each factor, which proves the first claim.
		\qedhere
	\end{itemize}
\end{proof}
Because it is not much harder, we also do the reductive case. By \Cref{cor:classify-semisimple-compact}, it is enough to reduce to the semisimple case.
\begin{proposition}
	Any connected compact Lie group $K$ is a finite quotient of a Lie group of the form $T\times C$, where $T$ is a product of $S^1$s, and $C$ is a simply connected semisimple compact Lie group.
\end{proposition}
\begin{proof}
	By \Cref{prop:compact-to-reductive}, $\mf k\coloneqq\op{Lie}K$ is reductive. Thus, $\mf k$ splits into a sum $\mf t\oplus\mf c$ where $\mf t$ is the center, and $\mf c$ is semisimple. It follows that the subgroup $T\subseteq K$ corresponding to $\mf t$ is central, and the subgroup $C\subseteq K$ corresponding to $\mf c$ is a normal subgroup. Note that $C$ is a closed subset in a compact space, so $C$ is compact.

	Now, the image of the homomorphism $T\times C\to K$ is a local diffeomorphism at the identity, so it contains an open neighborhood of the identity, so it is surjective. Thus, the kernel of this map is some discrete subgroup $\Gamma\subseteq T\times C$, so by compactness, $\Gamma$ is finite, and the result follows!
\end{proof}
\begin{remark}
	The proof also shows that we can take $T$ to be taken to be the connected component of the center of $K$.
\end{remark}

\subsection{Polar Decomposition}
As an aside to our study of real forms, we note that we are able to generalize the decompositions
\[\op{GL}_n(\RR)=\op B_n(\RR)\op O_n(\RR),\]
where $\op B_n(\RR)$ is the subgroup of upper-triangular matrices.

To this end, we need to generalize what we mean by $\op B_n$ and $\op O_n$ for general real forms. Let's start with the group.
\begin{notation}
	Fix a complex semisimple Lie algebra $\mf g$. For a Cartan involution $\theta$ inducing a real form $\mf g_\theta$, we let $G_{\mathrm{ad},\theta}\subseteq G_{\mathrm{ad}}$ be the corresponding subgroup, which are the fixed points of $\omega_\theta\coloneqq\omega\circ\theta$.
\end{notation}
\begin{example}
	The group $G_{\mathrm{ad},\theta}$ is not aways connected. For example, for $\mf g=\mf{sl}_2(\CC)$, take $\theta$ to induce the split form. Then $G_{\mathrm{ad},\theta}=\op{PSL}_2(\RR)$, which has two connected components: one can have positive or negative determinant!
\end{example}
Here is our theorem.
\begin{theorem}[Polar decomposition] \label{thm:polar-decomp}
	Fix a complex semisimple Lie algebra $\mf g$. For a Cartan involution $\theta$ inducing a real form $\mf g_\theta$. Let $K^c_\theta\subseteq G_{\mathrm{ad},\theta}$ be the subgroup of elements which act on $\mf g$ by unitary operators, and let $P_\theta$ be the image $\exp(i\mf p^c)$. Then the multiplication map
	\[\mu\colon K^c\times P_\theta\to G_{\mathrm{ad},\theta}\]
	is a diffeomorphism.
\end{theorem}
\begin{remark}
	By construction, we see that $P_\theta$ is the subgroup of $G_{\mathrm{ad},\theta}$ which acts on $\mf g$ by positive Hermitian operators. Thus, there is a $\log\colon P_\theta\to i\mf p^c$, which makes $P_\theta$ into a manifold.
\end{remark}
\begin{proof}
	The main point is to define an inverse for the multiplication map. We use the polar decomposition for usual matrices: any invertible complex matrix $A$ can be written uniquely as $A=U_AR_A$, where $U_A$ is unitary and $R_A$ is positive; explicitly, one can define $R_A\coloneqq\sqrt{A^\dagger A}$ and then define $U_A\coloneqq AR_A^{-1}$.

	Now, we may view $g\in G_{\mathrm{ad},\theta}$ as elements of $\op{Aut}(\mf g)$. For example, $g^\dagger g\in\op{Aut}(\mf g)$ has positive eigenvalues, and it lies in $G_{\mathrm{ad},\theta}$ because this element commutes with $\omega_\theta$. Thus, this element lives in $P_\theta$, and by undoing the exponential, we may take the square root to give $\sqrt{g^\dagger g}$. Thus, we may define $U_g$ and $R_g$ in $G_{\mathrm{ad},\theta}$ using the usual polar decomposition. By construction, $R_g\in P_\theta$, and we can see that $U_g$ acts by a Hermitian operator on $\mf g$, so $U_g\in K^c$ as well. Thus, $g\mapsto(U_g,R_g)$ is the required inverse for the multiplication map; note that it is smooth by construction!
\end{proof}
\begin{remark}
	It follows that $G_{\mathrm{ad},\theta}$ is homotopy-equivalent to $K^c$. Indeed, $P_\theta$ is diffeomorphic to a Euclidean space by construction.
\end{remark}
Here is an application.
\begin{corollary}
	Fix a semisimple complex Lie group $G$.
	\begin{listalph}
		\item The center $Z$ of $G$ is contained in $G^c$.
		\item Restriction defines an equivalence $\op{Rep}_\CC G\to\op{Rep}_\CC G^c$.
	\end{listalph}
\end{corollary}
\begin{proof}
	Let's start with (a); (b) then follows by using that the center is the obstruction to the descent of a representation from the simply connected cover, as in \Cref{cor:classify-semisimple-compact}. By \Cref{thm:polar-decomp}, we see that $G_{\mathrm{ad}}$ and $G_{\mathrm{ad}}^c$ are homotopy-equivalent, so their fundamental groups are the same. This fundamental group embeds into $G^c_{\mathrm{ad}}$ as some finite discrete central subgroup, where we see it must also embed into $G_{\mathrm{ad}}$. The case of general $G$ then follows by the classification given by \Cref{cor:classify-semisimple-compact}.
\end{proof}
\begin{remark}
	The same sort of arguments shows that there is also a polar decomposition to the real form $G_\theta$ of any connected complex semisimple group $G$.
\end{remark}
\begin{example}
	There is a notable caveat that even if $G$ is simply connected, the real form $G_\theta$ need not be connected and can even have infinite fundamental group. Indeed, for $G=\op{SL}_2(\CC)$, the split real form $\op{SL}_2(\RR)$ has infinite fundamental group: the polar decomposition gives
	\[\op{SL}_2(\RR)\cong\op{SO}_2(\RR)\times\RR^2,\]
	and $\op{SO}_2(\RR)\cong S^1$ has infinite fundamental group. In general, when $\mf k=\mf g^\theta$ fails to be semisimple, one finds that $G_\theta$ has some torus component, which gives the infinite fundamental group.
\end{example}
\begin{remark}
	When $G_\theta^\circ$ has infinite fundamental group, it turns out that the nontrivial covers of $G_\theta^\circ$ fail to be linear groups. As another example, for the split form $G_\theta=\op{Sp}_{2n}(\RR)$, one has $\mf k_\theta=\mf u(n)$, so $G_\theta$ is homotopy-equivalent to $\op U(n)$, which has fundamental group $\ZZ$. And indeed, the ``generalized metaplectic'' covers of $G_\theta$ fail to be linear! Roughly speaking, the issue is that $\op{Sp}_{2n}(\CC)$ has finite fundamental group, so any representation of $\widetilde{\op{Sp}}_{2n}(\RR)$ descends to the Lie algebra and then lifts to a finite cover of $\op{Sp}_{2n}(\CC)$, so it cannot be faithful.
\end{remark}

\end{document}